\chapter{Catalog, library and graph}
\label{ch:catalog}

Most organizations already have a report register --- an Excel sheet listing
every report, its vehicle, its discipline and where the file lives on a share.
The catalog features connect that register to the ingested corpus, so a question
can be answered about \emph{reports that exist} and not only about reports that
happen to have been ingested.

\section{The catalog}

\subsection{Import}

\route{POST /catalog/import} accepts \file{.xlsx} (via \code{openpyxl}) or a
delimited text file. Delimiter detection for text is automatic: tab, semicolon
or comma, whichever the first 2\,KB suggests.

A row must have at least six columns, and the first four must be non-empty.

\begin{longtable}{L{0.10\textwidth} L{0.22\textwidth} L{0.56\textwidth}}
\toprule
\textbf{Column} & \textbf{Field} & \textbf{Note} \\
\midrule
\endfirsthead
\bottomrule
\endfoot
1 & \code{report\_code} & Required. Also inspected as a possible file path. \\
2 & \code{vehicle\_name} & Required. \\
3 & \code{report\_title} & Required. \\
4 & \code{discipline} & Required. Upper-cased on import. \\
5 & \code{report\_date} & Optional. Dates are stored as ISO strings. \\
6 & \code{authors} & Optional. \\
\end{longtable}

Header rows are detected by their wording and skipped. A cell hyperlink is the
preferred source of the file path --- an Excel register usually links the report
rather than spelling the path out; \code{file:} URLs are decoded,
\code{http(s)} links are ignored, and separators are normalized to Windows form.

Each row is hashed. Re-importing the same sheet creates nothing new; if a row now
carries a source path it did not have before, only that path is updated. The
response reports \code{rows\_seen}, \code{created\_count},
\code{duplicate\_count}, \code{updated\_count} and the first twenty row-level
errors.

\subsection{Searching the catalog}

\route{GET /catalog/search} filters by free text, vehicle and discipline.
\route{GET /catalog/table} returns the register as a table, joined with what is
known about each row's ingestion state --- whether a document is linked, how many
chunks it has, and whether those chunks are embedded.

\subsection{Catalog questions}

\route{POST /ask/catalog} answers questions about the \emph{register itself}
rather than report content. \code{CatalogService} profiles the question (vehicles
mentioned, discipline, year, report code, intent) and then answers one of:

\begin{itemize}
\item an \textbf{analysis-type summary} --- what kinds of analysis exist for a
vehicle;
\item a \textbf{vehicle comparison} --- two or more vehicles side by side, with a
declared winner where the question asks for one; or
\item a \textbf{ranking} --- vehicles ordered by report count within a discipline
or year.
\end{itemize}

Everything is computed from catalog rows with SQL and deterministic profiling; no
model is involved.

\subsection{Linking catalog rows to files}

The link between a register row and a real file is the interesting part, because
the file lives on a share whose layout nobody controls.

\code{CatalogIngestService} searches the roots in \env{CATALOG\_SEARCH\_ROOTS}
with hard bounds --- maximum directory visits, maximum depth, maximum files per
directory, and a wall-clock deadline on the whole share scan. A disconnected
network drive slows a request down; it cannot pin a worker indefinitely.

\begin{longtable}{L{0.44\textwidth} L{0.46\textwidth}}
\toprule
\textbf{Endpoint} & \textbf{Purpose} \\
\midrule
\endfirsthead
\toprule
\textbf{Endpoint} & \textbf{Purpose} \\
\midrule
\endhead
\bottomrule
\endfoot
\route{GET /catalog/\{id\}/file-candidates} & The files that might be this row's
report, scored. \\
\route{GET /catalog/\{id\}/file-preview} & Preview one specific candidate. \\
\route{GET /catalog/\{id\}/best-file-preview} & Preview the best candidate. \\
\route{GET /catalog/\{id\}/best-file-preview-info} & Metadata about that
candidate without fetching it. \\
\route{POST /catalog/\{id\}/open-best-file} & Open it in the operating system's
default application (server-side, Windows). \\
\route{POST /catalog/ingest-candidate} & Ingest one chosen candidate and link it
to the row. \\
\route{POST /catalog/ingest-sample} & Background job: ingest a few reports per
discipline. \code{dry\_run} defaults to \code{true}, so the safe call is the
default one. \\
\route{POST /catalog/ingest-selected} & Background job: ingest the listed catalog
rows. \\
\route{POST /catalog/reconcile-documents} & Re-check existing links against the
documents in the database; supports \code{dry\_run}. \\
\end{longtable}

Links are stored in \code{catalog\_document\_links}, at most one document per
catalog entry, with the \code{match\_method} that produced the link.

\section{The library browser}

\route{POST /library/scan} returns a bounded, read-only tree of the supported
report files under a path. It exists for the \code{repocto} variant's file
browser, and it is written defensively:

\begin{itemize}
\item the path must sit inside an allowed root (\env{REPOCTO\_LIBRARY\_ROOTS}
plus the built-in defaults) --- anything else is refused;
\item the walk is capped at 800 documents, 800 directories, depth 10, and 2000
entries per directory; and
\item only \file{.pdf}, \file{.docx} and \file{.pptx} appear in the result.
\end{itemize}

It never reads file contents, and it never ingests anything.

\section{The graph overview}

\route{GET /graph/overview} builds a small graph of the corpus from catalog
entries: nodes for vehicles, disciplines, years and authors, with edges to the
reports that connect them, plus tag counts. The result is bounded (\code{limit}
defaults to 160, clamped to 20--300) because the view is meant to be readable,
not exhaustive.

The \code{ANALYSIS TYPE} placeholder discipline --- a header artefact that
survives some registers --- is filtered out here, as it is everywhere else in the
catalog code.
